%! Inverse Functions — Unit 1, Lesson 1.6 — Lesson Plan
\documentclass[10pt]{article}
\usepackage{precalculus-article}
\usepackage{precalculus-boxes}
\usepackage{graphicx}
\graphicspath{{images/}}
\newcommand{\TallMath}[1]{$\displaystyle #1\rule[-1.4em]{0pt}{3.2em}$}
\newcommand{\CourseName}{Precalculus}
\newcommand{\MeetingLength}{55 minutes}
\newcommand{\UnitNumberName}{Unit 1: Functions and Change \quad}
\newcommand{\LessonNumberName}{Lesson 1.6: Inverse Functions}

\begin{document}
\begin{center}
    {\Large\bfseries \CourseName \\
    {\small\bfseries \UnitNumberName \LessonNumberName}}
\end{center}

\begin{tcolorbox}[colback=lilac, colframe=plum, boxrule=0.9pt, arc=2mm,
    left=2mm, right=2mm, top=1.5mm, bottom=1.5mm]
    \textbf{Primary Objective:} Students will build the inverse of a function by
    reversing its input--output pairs and by the analytic swap-and-solve method;
    will verify a pair of inverses by composition, using
    $f(f^{-1}(x)) = f^{-1}(f(x)) = x$; and will decide when a domain must be
    restricted before a function can be inverted.
    \hfill\textit{\small (CED Topic 2.8 --- Skills 1.A, 2.B)}
\end{tcolorbox}

\renewcommand{\arraystretch}{1.6}
\begin{skillbox}[Priority Ideas \& Skills]{goldbox}
\begin{tabularx}{\linewidth}{@{}X X@{}}
\textbf{AP Skills} &
\textbf{Key Understandings} \\
\midrule
\textbf{1.A} --- Solve equations and inequalities represented analytically,
with and without technology. &
\begin{itemize}[leftmargin=1.2em,itemsep=2pt,topsep=0pt]
  \item A function is invertible on a domain exactly when each output comes
        from a unique input --- it is \textbf{one-to-one} there. Otherwise the
        domain must be \textbf{restricted}. (2.8.A.1)
  \item The analytic method: write $y = f(x)$, swap $x$ and $y$, solve for $y$,
        and rename the result $f^{-1}(x)$. (2.8.B.4)
\end{itemize} \\
\textbf{2.B} --- Construct equivalent graphical, numerical, analytical, and
verbal representations of functions. &
\begin{itemize}[leftmargin=1.2em,itemsep=2pt,topsep=0pt]
  \item If $f(a) = b$ then $f^{-1}(b) = a$: the inverse reverses every
        input--output pair, so the \textbf{domain and range swap}. (2.8.A.2)
  \item The graph of $y = f^{-1}(x)$ is the reflection of $y = f(x)$ over the
        line $y = x$. (2.8.B.3)
\end{itemize} \\
\textbf{Synthesis} --- inverses are checked, not asserted. &
\begin{itemize}[leftmargin=1.2em,itemsep=2pt,topsep=0pt]
  \item Verification by composition: $f(f^{-1}(x)) = f^{-1}(f(x)) = x$, the
        identity function. Lesson 1.5 built exactly this tool. (2.8.B.1)
  \item \textbf{Notation warning:} $f^{-1}(x) \ne \dfrac{1}{f(x)}$. The $-1$ is
        a name, not an exponent --- today's signature error.
\end{itemize} \\
\end{tabularx}
\end{skillbox}

\begin{skillbox}[{Vocabulary, Concepts, \& Theorems}]{greenbox}
\renewcommand{\arraystretch}{1.4}
\begin{tabularx}{\linewidth}{@{} >{\leavevmode\bfseries\color{plum}}p{0.28\linewidth} X @{}}
Inverse function      & The function $f^{-1}$ that reverses $f$: if $f(a) = b$, then $f^{-1}(b) = a$. Read ``$f$-inverse.'' \\
One-to-one            & No two different inputs produce the same output. A function is invertible on a domain exactly when it is one-to-one there. \\
Domain restriction    & Cutting a function's domain down to a piece on which it \textit{is} one-to-one, so an inverse exists there. \\
Identity function     & The function that returns its own input, $I(x) = x$. Composing a function with its inverse produces it, in either order. \\
Reflection over $y=x$ & The graph of $y = f^{-1}(x)$ is the mirror image of $y = f(x)$ across the line $y = x$; every point $(a,b)$ becomes $(b,a)$. \\
$f^{-1}$ vs.\ $1/f$   & \TallMath{f^{-1}(x) \ne \frac{1}{f(x)}} --- the $-1$ names the reversing function; it is not a reciprocal or an exponent. \\
\end{tabularx}
\end{skillbox}

\begin{skillbox}[Activate Prior Knowledge \& Spiral Review]{lilac}
    Please see attached warm-up (spiral review).
    Skills reviewed: solving a linear equation for $x$ (the algebra that Step 3
    of the analytic method reduces to), evaluating a composition numerically
    from Lesson 1.5, and reversing a small table of input--output pairs.
    Item 2 is the load-bearing one --- students compute $f(g(7)) = 7$ and
    $g(f(7)) = 7$ without being told why, and Part 4 reveals that they were
    verifying a pair of inverses the whole time.
\end{skillbox}

\begin{skillbox}[Hook]{lilac}
    \textbf{Give me back the Celsius.} Project the weather app: $68^\circ$F.
    The only conversion rule on the board is $F(x) = 1.8x + 32$, which turns
    Celsius \textit{into} Fahrenheit --- the wrong direction.
    \begin{enumerate}[itemsep=2pt]
        \item Ask: what do you do to $68$ to get back to Celsius, and in what
              \textit{order}? (Subtract 32 first, then divide by 1.8 --- undo
              the last step first.)
        \item Then press: you just followed a rule. Write it down. That new
              rule is a function too, and it has a name.
    \end{enumerate}
    \textit{Discussion prompt:} every function that never repeats an output has
    a partner that runs it backwards. Today we build the partner, check it, and
    learn when it refuses to exist.
\end{skillbox}

\begin{skillbox}[Lesson]{lilac}
\begin{multicols}{2}

\textbf{Part 1: ``Undoing'' as a Reversed Mapping} (8 min)

\begin{itemize}
  \item State the definition once, in the form students will use all lesson:
        \textbf{if $f(a) = b$, then $f^{-1}(b) = a$} (2.8.A.2). Read it aloud
        both ways off the temperature rule: $F(20) = 68$, so
        $F^{-1}(68) = 20$.
  \item Deliver the notation warning \textit{now}, before any algebra:
        $f^{-1}(x)$ is \textbf{not} $1/f(x)$. Make it numerical rather than a
        slogan --- $F(10) = 50$ gives $F^{-1}(50) = 10$, while
        $1/F(50) = 1/122$. Two wildly different numbers.
\end{itemize}

\textbf{Part 2: Inverses from Tables --- the Domain/Range Swap} (10 min)

\begin{itemize}
  \item Fill the $F$ table, then build the $F^{-1}$ table by trading the two
        rows. Nothing is computed; the pairs are re-read backwards.
  \item Name the consequence explicitly: the \textbf{domain of $f^{-1}$ is the
        range of $f$}, and vice versa (2.8.A.2).
  \item Counterexample: $g = \{(1,5),(2,5),(3,8)\}$. Reversed, the input 5
        would need two outputs --- so the reversal is not a function
        (2.8.A.1).
\end{itemize}

\textbf{Part 3: Swap and Solve} (14 min)

\begin{itemize}
  \item The four steps: write $y = f(x)$; \textbf{swap} $x$ and $y$; solve for
        $y$; rename $f^{-1}(x)$ (2.8.B.4). Step 2 is the inversion --- steps 3
        and 4 are bookkeeping.
  \item Do it twice, deliberately. First $f(x) = 4x - 1$, giving
        $f^{-1}(x) = \dfrac{x+1}{4}$; then the lesson's context,
        $F(x) = 1.8x + 32$, giving $F^{-1}(x) = \dfrac{x-32}{1.8}$.
  \item Close the loop: evaluate $F^{-1}(68)$ and get $20$ --- the same number
        the table gave in Part 2.
\end{itemize}

\columnbreak

\textbf{Part 4: Verification by Composition} (10 min)

\begin{itemize}
  \item \textbf{This is where Lesson 1.5 pays off.} Composition was defined
        there; here it becomes a test. Two functions are inverses exactly when
        \textbf{both} compositions collapse to the identity:
        $f(f^{-1}(x)) = f^{-1}(f(x)) = x$ (2.8.B.1).
  \item Verify $f(x) = 4x-1$ against $f^{-1}(x) = \dfrac{x+1}{4}$ in both
        orders, line by line. Then point back at the warm-up: they already got
        $7$ both ways with these exact functions. Algebra just does every input
        at once.
  \item Insist on \textit{both} directions. One composition alone does not
        settle it.
\end{itemize}

\textbf{Part 5: No Inverse --- and the Mirror Line} (8 min)

\begin{itemize}
  \item $f(x) = x^2$ fails: $f(-2) = f(2) = 4$, so reversing sends 4 to two
        places. Restrict to $x \ge 0$; then $f^{-1}(x) = \sqrt{x}$ (2.8.A.1).
  \item Read the pre-drawn figure: $y = f^{-1}(x)$ is the reflection of
        $y = f(x)$ over $y = x$, and $(0,1)$ on one graph corresponds to
        $(1,0)$ on the other (2.8.B.3). Students read the figure --- they do
        not draw it.
\end{itemize}
\end{multicols}
\end{skillbox}

\begin{skillbox}[Active Monitoring]{redbox}
While students work on guided notes and practice, circulate and check:
\begin{itemize}
  \item \textbf{Common error:} writing $f^{-1}(x) = \dfrac{1}{f(x)}$. This is
        the lesson's signature error. Cold-call with numbers, never with a
        rule: ``$F(10) = 50$. What is $F^{-1}(50)$? Now what is $1/F(50)$?
        Are those the same number?''
  \item \textbf{Common error:} skipping the swap in Step 2 and simply solving
        the original equation for $x$. Ask: ``Which variable is the input of
        the \textit{new} function?''
  \item \textbf{Common error:} in the table reversal, recomputing outputs
        instead of trading rows. Nothing is calculated in Part 2 --- the same
        pairs are read backwards.
  \item \textbf{Common error:} verifying only $f(f^{-1}(x))$ and stopping.
        Both directions, every time.
  \item \textbf{Common error:} claiming $f(x)=x^2$ ``has no inverse.'' It has
        none \textit{on all real numbers}; on $x \ge 0$ it has one. Push for
        the domain in the sentence.
\end{itemize}
\end{skillbox}

\begin{skillbox}[Group Work \& Differentiation]{redbox}
Students work on the \textbf{Group Activity: Reading the Thermometer
Backwards} in leveled tiers. Every tier uses the same conversion rule
$F(x) = 1.8x + 32$ and the same table --- no tier draws anything.

\begin{itemize}
  \item \textbf{Tier R (Remediate):} Complete the Fahrenheit row, build the
        $F^{-1}$ table by trading rows, read $F^{-1}(68)$ off it, and state
        the domain and range of both functions to see them swap.
  \item \textbf{Tier A (Apply):} Find $F^{-1}(x)$ by swap-and-solve, verify
        with $F(F^{-1}(x))$, evaluate $F^{-1}(212)$, and write one sentence
        interpreting what $F^{-1}$ does in context.
  \item \textbf{Tier E (Extend):} A greenhouse temperature
        $T(t) = (t-6)^2 + 50$ is not one-to-one on $0 \le t \le 12$ --- find
        two times with the same temperature, restrict the domain, invert on
        the restriction, and finish with the $F^{-1}(50)$ versus $1/F(50)$
        comparison.
\end{itemize}
\end{skillbox}

\begin{skillbox}[Individual Work \& Assessment]{redbox}
\textbf{Exit Ticket} (last 5--7 minutes, individual, collected):
\begin{enumerate}
  \item Given $f = \{(1,4),(2,9),(3,16)\}$, write $f^{-1}$ as ordered pairs and
        state its domain.
  \item Find $f^{-1}(x)$ for $f(x) = 2x - 5$ (worked space provided).
  \item Verify by computing $f(f^{-1}(x))$, then circle the true statement
        about $f^{-1}$ versus $1/f$.
\end{enumerate}
\textbf{Assessment note:} Item 1 checks the reversal idea (2.8.A); items 2--3
check the analytic method and verification (2.8.B). The circle in item 3 is the
diagnostic --- a student who circles $f^{-1}(x) = 1/f(x)$ needs Part 1 retaught
in tomorrow's warm-up window regardless of how items 1 and 2 went.
\end{skillbox}

\begin{skillbox}[Reinforcement \& Extension]{goldbox}
\textbf{Homework:} 6 problems covering table reversal and invertibility, two
swap-and-solve inverses, verification in both directions, the temperature
context (including the $f^{-1}$ versus $1/f$ comparison), domain restriction
for $f(x) = x^2$, and one AP-style multiple-choice item on what $f(3) = -2$
tells you about $f^{-1}$.

\textbf{Extension (optional):} Two functions are given whose compositions are
\textit{almost} the identity --- $f(f^{-1}(x)) = x$ but the second composition
is not $x$ on the stated domain. Students find where the claim breaks and
repair it with a domain restriction.

\textbf{Preview:} Next lesson (1.7, Function Model Selection and Construction)
asks which kind of function fits a context --- and inverting a model is how you
answer ``what input produced this output?'' Long-range: in Unit 5 the
logarithm arrives as nothing more than the inverse of the exponential
function, built by exactly today's method.
\end{skillbox}

\begin{teachernote}[Warm-Up]
Five minutes, silent start, no calculators except for item 1(b). Item 2 is the
one that pays off: students compute $f(g(7))$ and $g(f(7))$ and both come back
to 7, with no explanation offered. Do not explain it at the warm-up --- let it
sit. In Part 4, put the warm-up back on the board and say ``you verified a pair
of inverses before I told you what one was.'' Item 3's table reversal is Part 2
in miniature; if a student recomputes outputs instead of trading rows, that is
the exact misconception Part 2 addresses.
\end{teachernote}

\begin{teachernote}[Guided Notes]
Front-load the notation warning in Part 1 and make it numerical --- ``$1/122$
versus $10$'' beats any amount of ``the $-1$ isn't an exponent.'' Part 3 is
where the time goes; do both worked examples, and do not let the class talk you
out of the second one because ``it's the same as the first.'' The repetition is
the point for this track. Part 4 is the payoff for Lesson 1.5: name the
identity function out loud, since 1.5's plan promised this lesson would deliver
it. If time runs short, compress Part 5's $x^2$ discussion to the single
sentence ``$f(-2)$ and $f(2)$ are both 4'' and keep the figure read --- the
reflection over $y = x$ is assessed on the homework.
\end{teachernote}

\begin{teachernote}[Group Activity]
All three tiers run on the same conversion rule and the same table, so groups
can be re-sorted mid-activity without new setup. Tier R's last question --- the
domain and range of $F$ and $F^{-1}$ side by side --- is that tier's whole
point; make them say the word ``swap'' out loud. Tier A's interpretation
sentence is worth more than its algebra: ``$F^{-1}$ takes a Fahrenheit reading
and returns the Celsius temperature'' is the target. Tier E's greenhouse
function is the only place a non-linear restriction is practiced; if the class
never reaches it, assign homework problem 5 as the substitute.
\end{teachernote}

\begin{teachernote}[Exit Ticket]
Collected, completion credit only. Sort by item 3's circled statement first,
then by item 2. Accept $\dfrac{x+5}{2}$ or $0.5x + 2.5$ on item 2 --- both are
correct and the second often signals a student who solved cleanly. On item 1, a
student who writes the pairs in a different order is fine; a student who writes
$\left\{\left(1,\tfrac14\right),\ldots\right\}$ has the reciprocal
misconception and goes in the reteach pile with the item-3 misses.
\end{teachernote}

\begin{teachernote}[Homework]
Problem 4 is the storyline problem and the one to scan first --- part (c) is
the notation warning in numbers, and a wrong answer there predicts trouble in
Unit 5. On problem 6, distractor A swaps the roles of input and output in the
inverse, C applies the reciprocal misconception to the output, and D applies it
to both --- each is worth thirty seconds of autopsy tomorrow. The extension is
genuinely hard for this track; assign it to students who finished Tier E, not
to everyone.
\end{teachernote}
\end{document}
